\chapter{任务三：Context Path、Provenance 与可验证证据}
\label{chap:task3}

\section{约束}

多轮 Agent 需要把工具结果反馈给下一轮，但“记住文本”不足以形成可信 Context。我们
必须区分谁产生了结果、它来自文件还是内核、通过哪条 route 跨 Agent、当前 active
branch 是什么，以及一次 rollback 是否只是改变后续视图而没有撤销既有外部效果。

此外，题面要求超过五轮并在超长时自动淘汰。固定容量必须同时避免 torn read、引用
已淘汰 predecessor 和把用户 cache 当成授权事实。我们用进程级 Context Path 保存
运行历史，用 workflow Evidence Ring 保存 canonical evidence，再由 fence 给出外部可
复核 cut。

\section{Context 记录与因果路径}

每次工具调用在效果边界形成 Context record。sequence 单调递增；记录绑定 request、
tool、status、service-start tick、cause sequence、span、branch generation、control id、
payload/result 摘要和 record hash。\code{path_parent_sequence} 指向当前 active head，
successor index 支持分支查询。

Context append 采用 prepare、填充、publish、commit 的顺序。只读 mirror header 同时
发布 count、head、oldest、latest、visible head、dropped 和 rollback count。读者先取
稳定 summary，再读 slot 并重验 sequence/hash，最后确认 summary 未变化，从而拒绝
torn snapshot。

容量为 128 条。写满后按 sequence 淘汰最旧记录；active path 不能引用已淘汰节点。
\code{agentfinal_ucore} 在一次 64-op compact batch 后继续执行 snapshot、branch、
rollback 和 window 淘汰，超过题面五轮门槛。

\section{Rollback 的精确语义}

rollback 只把 \code{visible_head} 移到仍保留的某个 Context sequence，并为后续调用
分配新 branch generation。它不删除历史 record，也不倒转已经发生的文件写入、IPC、
进程创建或外部 provider 请求。clear 同样重置当前 Context 状态，而不是系统级事务
回滚。

这个限制非常重要。如果把 Context rollback 描述成副作用回滚，用户会错误相信一次
已批准动作可以被撤销。真正需要效果幂等或补偿时，应由 V3 attempt cache、工具协议或
业务任务显式实现。

\section{六类固定 provenance 标签}

我们不让内核解析自由文本或运行 prompt-injection 分类器，而是使用六个可组合标签：

\begin{table}[htbp]
\centering
\caption{Context provenance 词汇}
\label{tab:provenance-labels}
\begin{tabularx}{\textwidth}{L{4.2cm}Y L{4.3cm}}
\toprule
\textbf{标签} & \textbf{含义} & \textbf{典型来源} \\
\midrule
\code{KERNEL_FACT} & 由内核结构化接口产生的事实 & PID、resource、scheduler、capability \\
\code{TRUSTED_USER_CONTROL} & 经过产品控制面定义的用户控制 & 新用户 turn、精确审批决定 \\
\code{AGENT_DERIVED} & Agent 从已有输入推导 & 摘要、报告、计划 \\
\code{UNTRUSTED_FILE_DATA} & 文件内容或元数据输入 & query/read/digest source \\
\code{UNTRUSTED_TOOL_OUTPUT} & 工具执行返回 & V2 result、模型工具结果 \\
\code{CROSS_AGENT_DATA} & 经过 IPC 或工件跨角色 & TASK result、broker artifact \\
\bottomrule
\end{tabularx}
\end{table}

文件、工具、IPC mailbox 和 Task resource 都执行保守 OR 传播。重新哈希、摘要或跨 Agent
materialize 不会删除旧标签。\code{send_message/llm_request/llm_response} 允许这些带
来源的观测继续进入 Agent Loop，但输出保留原污点，route、capability、schema 和 effect
检查不放松。其他副作用工具不会因为数据能被发送而获得额外权限。

\section{Manifest 与 pre-effect gate}

每个内核工具的安全 manifest 声明 accepted input labels、output-added labels、required
capability 以及 file/metadata/IPC/process/permission/artifact/watch side-effect mask。
启用 ENFORCE V3 时，外部效果必须同时满足 lifecycle generation、frozen contract edge、
schema、capability、manifest 和 provenance。任何一项不满足都在效果前返回
\code{DENIED}，并把 source Context、tool、reason、lifecycle 和 evidence ticket 写入
critical evidence。

V1/V2 仍检查 capability、scope 与 schema，但没有 frozen-edge/provenance envelope。
我们在文档和产品中保持这一区分，避免把 Nexus V2 的灵活性包装成 V3 的高保证。

\section{Fence-Sealed Evidence Ring}

每个活跃 workflow 按需分配四个 Agent state page：三页 ordinary ring，共 48 个
256-byte slot；一页 critical ring，共 16 个 slot。生产者 reserve 单调 ticket，在
IRQ-on 区域填充后 commit，失败则 discard 并记录 gap。读者遇到更早 BUSY ticket 时
隐藏后续记录，保持发布顺序。

普通成功 Context 只写一次 canonical evidence，避免同一事实同时写 Context、ledger
和多个 hash chain。拒绝或授权效果进入 critical ring，并保留一条兼容 ledger 投影；
ring 失败则 fail closed 到 legacy ledger。audit/timeline/provenance 因而仍是合并视图，
不是已经删除的旧接口。

环满时，已提交 segment 先滚入内部 root 再复用槽。显式 workflow fence 使用 SHA-256
domain separation，绑定 previous public root、challenge、fence sequence、ticket range、
event/gap count、metadata generation、credit epoch、exact U digest 和 ordered segment
digest。只有 challenge-bound fence 才推进外部 root；内部 rollover seal 不冒充 receipt。

\section{Workflow fence 的精确 cut}

controller 通过 \code{agent_run(count=0, AGENT_RUN_F_FENCE)} 发起 fence。内核依次检查
orchestrator/control、关闭 gate、取得 metadata quiescence、排空 live-query deferred
work、处理 VFS deferred reclaim、trim exec/storage credit 并要求所有 P 为零，随后
密封 Evidence Ring。receipt 在内部状态提交后才 copyout；copyout 失败的同 id 重试返回
缓存结果，不重新切割。

320-byte receipt 明确标记 partial coverage、credit exact、evidence sealed 和 metadata
volatile。Evidence Ring 是内存结构；sealed 不表示落盘、fsync、掉电原子或跨重启恢复，
也不覆盖全部文件内容、Host 行为和每条调度事件。

\section{验证}

\code{agentfinal_ucore} 验证 64 次顺序调用、可信 mirror、query/snapshot、branch、rollback、
淘汰和 record hash。provenance fixture 在 ENFORCE V3 下把不可信来源导向计划外副作用，
要求 pre-effect \code{DENIED} 和 critical evidence；mutation tests 删除标签传播或 effect
mask 时必须失败。fence tests 覆盖 challenge、retry cache、conflict/stale、credit P、
metadata generation 和 copyout cut。

\begin{evidencebox}[任务三结论]
Context 是有界、可查询、带因果和来源的内核历史；Evidence Ring 与 fence 提供当前 boot
的 challenge-bound cut。它不证明文本语义安全，也不是持久日志或远程证明系统。
\end{evidencebox}
